\documentclass[12pt]{article}
\usepackage[utf8]{inputenc}
\usepackage[T1]{fontenc}
\usepackage{hyperref}
\title{Pulaar Language Resources for K-12 Curriculum Development}
\author{Compiled by Claude Sonnet 5.5 (claude-sonnet-5-5) for Andr\'as Kornai}
\date{}
\begin{document}
\maketitle
\section{Introduction}
Pulaar is spoken in Senegal, Mauritania, Guinea, The Gambia, Guinea-Bissau and Mali. The search found mainly English-language learner and reference materials; native-language school materials were not found in results.
\section{Primary Educational Resources}
\begin{itemize}
\item Peace Corps Pulaar (Mauritania) language course, a commercial listing: \url{https://www.fsi-language-courses.org/pulaar/courses/peace-corps-pulaar-mauritania-language-course}
\item Columbia University Fula library guide (dictionaries, grammars, readers, textbooks): \url{https://guides.library.columbia.edu/fula-language/home}
\end{itemize}
\section{Linguistic and Reference Materials}
\begin{itemize}
\item Pulaar--English/English--Pulaar Standard Dictionary, about 30,000 entries (Hippocrene): \url{https://www.hippocrenebooks.com/store/p275/Pulaar-English_StandardDictionary.html}
\item Fula dictionary links (Lexilogos): \url{https://www.lexilogos.com/peul_dictionnaire.htm}
\item MSU ALMA bibliography listing a 1995 study of reading and education in Pulaar in Senegal and a 1993 Pulaar-Pulaar-French thematic dictionary (Dakar): \url{https://alma.matrix.msu.edu/assets/uploads/2009/04/bibliographyfarysilateka.pdf}
\end{itemize}
\section{Religious and Community Resources}
None identified in search results.
\section{Digital and Audio Resources}
None identified beyond the library guide.
\section{Curriculum Development Context}
Pulaar is widely spoken, but no national school curricula were found. Ministries of education and SIL/ALMA archives should be consulted next. The Pulaar Learners' Reference Grammar (2002) was mentioned in a search summary without a URL, so it is not listed.
\end{document}